\section{Binary Expand--Convolute}
\label{sec:binary-ec}

The accumulator has one state bit.  A random convolution has $m$ state bits,
but its output-weight law does not require all $2^m$ bit patterns.  It is
enough to record how many consecutive zeros end the current output prefix.
This reduction gives an exact transfer matrix with $m+1$ states.

We analyze the wrapped convolution from \eqref{eq:wrapped-convolution}.  The
fixed oldest tap forces the process to leave the almost-zero state unless the
next input cancels it.  This behavior improves the sparse-message exponent.

\subsection{Exact transfer matrix}

Index the reduced states by $0,\ldots,m$.  State $j<m$ means that the current
output prefix ends in exactly $j$ zeros.  State $m$ means that its last $m$
outputs are all zero.  The initial state is $m$.

Let $X$ mark nonzero input bits and let $Y$ mark nonzero output bits.  Define
the $(m+1)\times(m+1)$ matrix $\mathbf W_m(X,Y)$ by its nonzero entries:
\begin{align}
 (\mathbf W_m)_{j,0}&=\frac{1+X}{2}Y,&
 (\mathbf W_m)_{j,j+1}&=\frac{1+X}{2}
 &&(0\leq j<m-1),
 \label{eq:ec-active-transitions}\\
 (\mathbf W_m)_{m-1,0}&=Y,&
 (\mathbf W_m)_{m-1,m}&=X,
 \label{eq:ec-critical-transitions}\\
 (\mathbf W_m)_{m,0}&=XY,&
 (\mathbf W_m)_{m,m}&=1.
 \label{eq:ec-zero-transitions}
\end{align}

\begin{lemma}[Exact wrapped-convolution enumerator]
\label{lem:ec-exact-enumerator}
Let $E^{\mathrm{wConv}}_{u,h}$ be the expected number of weight-$u$ inputs
that a sampled wrapped convolution maps to weight $h$.  Then
\begin{equation}
 E^{\mathrm{wConv}}_{u,h}=
 [X^uY^h]\,\e_m^{\mathsf T}\mathbf W_m(X,Y)^n\one,
 \label{eq:ec-exact-enumerator}
\end{equation}
where $\e_m$ selects the all-zero initial state.
\end{lemma}

\begin{proof}
If $j<m-1$, at least one state bit with a random tap is nonzero.  Its inner
product with the sampled taps is uniform, independently of the input bit.
This gives \eqref{eq:ec-active-transitions}.  In state $m-1$, only the fixed
oldest tap is active.  Input zero produces output one; input one produces zero
and enters state $m$.  This gives \eqref{eq:ec-critical-transitions}.  In
state $m$, the feedback is zero and the output equals the input.  Matrix
multiplication concatenates the transitions.
\end{proof}

If the input bits are independent $\Ber(a)$ variables, averaging the input
marker gives a one-variable matrix $\mathbf N_{m,a}(Y)$ with nonzero entries
\begin{align}
 (\mathbf N_{m,a})_{j,0}&=Y/2,&
 (\mathbf N_{m,a})_{j,j+1}&=1/2
 &&(0\leq j<m-1),
 \label{eq:ec-ber-active}\\
 (\mathbf N_{m,a})_{m-1,0}&=(1-a)Y,&
 (\mathbf N_{m,a})_{m-1,m}&=a,
 \label{eq:ec-ber-critical}\\
 (\mathbf N_{m,a})_{m,0}&=aY,&
 (\mathbf N_{m,a})_{m,m}&=1-a.
 \label{eq:ec-ber-zero}
\end{align}

\begin{theorem}[Exact Bernoulli wrapped-EC first moment]
\label{thm:ec-exact-first-moment}
Independently sample
$\B\gets\mathcal B^{\mathrm{Ber}}_{k,n,\rho}$ and a wrapped convolution
$\C$ of memory $m$.  Let $a_r$ be defined by \eqref{eq:ea-activation}.  Then
\begin{equation}
 \Pr_{\B,\C}[\Bad_L(\B\C)]
 \leq
 \sum_{r=1}^{k}\binom{k}{r}
 \sum_{h=0}^{L}[Y^h]\,
 \e_m^{\mathsf T}\mathbf N_{m,a_r}(Y)^n\one.
 \label{eq:ec-first-moment}
\end{equation}
\end{theorem}

\begin{proof}
For a fixed weight-$r$ message, the Bernoulli expander produces independent
$\Ber(a_r)$ inputs to the convolution.  Matrix $\mathbf N_{m,a_r}$ gives the
exact output-weight generating function.  Sum over output weights and
messages, then apply \eqref{eq:first-moment}.
\end{proof}

For a numerical marker $z\in(0,1]$, the inner tail in
\eqref{eq:ec-first-moment} is at most
\begin{equation}
 z^{-L}\e_m^{\mathsf T}\mathbf N_{m,a_r}(z)^n\one.
 \label{eq:ec-marker-bound}
\end{equation}
The calculation uses $m+1$ states and remains exact before the positive-marker
bound.

\subsection{Asymptotic distance}

Let $\varrho_m(a,z)$ be the spectral radius of
$\mathbf N_{m,a}(z)$, namely the largest absolute value of its eigenvalues.
Define
\begin{align}
 J^{\mathrm w}_{m,\delta}(a)
 &:=\sup_{0<z\leq1}
 \{\delta\ln z-\ln\varrho_m(a,z)\},
 \label{eq:ec-rate-function}\\
 I^{\mathrm w}_{m,\delta}
 &:=\left(
 \sqrt{1-2\delta+2^{1-m}\delta}
 -\sqrt{2^{1-m}\delta}
 \right)^2.
 \label{eq:ec-sparse-rate}
\end{align}

\begin{lemma}[Wrapped-convolution endpoint exponents]
\label{lem:ec-endpoints}
For fixed $m\geq1$ and $\delta\in(0,1/2)$,
$J^{\mathrm w}_{m,\delta}$ is continuous and positive on $(0,1/2]$.
Moreover,
\begin{align}
 \liminf_{a\downarrow0}
 \frac{J^{\mathrm w}_{m,\delta}(a)}a
 &\geq I^{\mathrm w}_{m,\delta},
 \label{eq:ec-sparse-endpoint}\\
 J^{\mathrm w}_{m,\delta}(1/2)
 &=\ln2-h(\delta).
 \label{eq:ec-dense-endpoint}
\end{align}
The associated lower-tail bounds can be chosen uniformly for $a$ near zero
and for $a$ in any closed interval bounded away from zero.
\end{lemma}

\begin{proof}
Perron--Frobenius theory bounds
$\e_m^{\mathsf T}\mathbf N_{m,a}(z)^n\one$ by a constant times
$\varrho_m(a,z)^n$.  At $a=0$ and $z=1$, the chain has an active class and an
all-zero class.  First-order perturbation with $z=e^{-\theta a}$ reduces the
calculation to a two-state continuous-time generator.  Optimizing its larger
eigenvalue gives \eqref{eq:ec-sparse-rate}.  At $a=1/2$, every row of
$\mathbf N_{m,1/2}(z)$ sums to $(1+z)/2$.  Hence
$\varrho_m(1/2,z)=(1+z)/2$, which gives
\eqref{eq:ec-dense-endpoint}.  Appendix~\ref{app:analytic-details} gives the
perturbation and proves a uniform quadratic remainder in
Lemma~\ref{lem:ec-two-class-perturbation}.  It also proves the uniform
prefactor claim.
\end{proof}

\begin{theorem}[Binary wrapped-EC distance]
\label{thm:ec-distance}
Fix $m\geq1$ and $\delta\in(0,1/2)$.  Let $k/n\to R\in(0,1)$ and set
$\rho_n=C\ln(n)/n$.  Independently sample
$\B\gets\mathcal B^{\mathrm{Ber}}_{k,n,\rho_n}$ and a wrapped convolution
$\C$ of memory $m$.  If
\begin{align}
 C I^{\mathrm w}_{m,\delta}&>1,
 \label{eq:ec-sparse-condition}\\
 h(\delta)&<(1-R)\ln2,
 \label{eq:ec-gv-condition}
\end{align}
then
\begin{equation}
 \Pr_{\B,\C}[d_{\min}(\B\C)\leq\lfloor\delta n\rfloor]=o(1).
 \label{eq:ec-distance-conclusion}
\end{equation}
\end{theorem}

\begin{proof}
Apply Lemma~\ref{lem:ec-endpoints} to the exact first moment in
Theorem~\ref{thm:ec-exact-first-moment}.  For sparse messages,
$a_r=r\rho_n(1-o(1))$, and \eqref{eq:ec-sparse-condition} makes the
support-$r$ contribution summable.  For intermediate weights, $a_r$ is
bounded away from zero, so a uniform exponential tail dominates the number
of messages.  For every $r>\eta n$, one has $a_r=1/2-o(1)$ uniformly.  The
dense endpoint and the bound of $2^k$ messages give exponent
\[
 R\ln2+h(\delta)-\ln2+o(1),
\]
which is negative by \eqref{eq:ec-gv-condition}.  These three ranges cover
$[1,k]$.
\end{proof}

For $m=1$, equation~\eqref{eq:ec-sparse-rate} equals the EA value
$I_\delta$.  As $m\to\infty$, it tends to $1-2\delta$.  At
$\delta=0.05$ and $m=21$, the new sparse condition is $C>1.111623$.
The published nonwrapping theorem requires $C>10.4344$ and $R<0.1278$ at the
same distance~\cite{EC}.  The wrapped theorem permits every $R<0.7136$, the
binary GV rate at distance $0.05$.

The asymptotic comparison changes the analyzed member of the EC family: the
published theorem is nonwrapping, while Theorem~\ref{thm:ec-distance} is for
the wrapped construction already proposed in~\cite{EC}.  The next comparison
holds the wrapped finite ensemble fixed.

\subsection{Finite same-ensemble comparison}

Set $k:=2^{20}$, $n:=5k$, $L:=n/20$, and $m:=5$.  For each row of
Table~\ref{tab:ec-same-ensemble}, set $\rho:=C_0\ln(n)/n$.  The boundary state
and expander distribution match the zero-state wrapped rows of
\cite[Table~4]{EC}.

\begin{theorem}[Finite wrapped-EC comparison]
\label{thm:ec-finite-comparison}
For each row, independently sample
$\B\gets\mathcal B^{\mathrm{Ber}}_{k,n,\rho}$ and a wrapped convolution
$\C$ of memory $5$.  The exact-transfer column of
Table~\ref{tab:ec-same-ensemble} upper-bounds
$\Pr_{\B,\C}[\Bad_L(\B\C)]$.
\end{theorem}

\begin{table}[t]
\centering
\caption{Same-ensemble comparison with the prior wrapped-EC calculation.}
\label{tab:ec-same-ensemble}
\begin{tabular}{@{}ccccc@{}}
\toprule
$C_0$ & expected degree & prior bound & exact transfer & failure exponent \\
\midrule
$3.0$ & $46.4171$ & $2.26\cdot10^{-7}$ & $2.187816572\cdot10^{-10}$ & $32.089789$ \\
$2.5$ & $38.6810$ & $1.72\cdot10^{-5}$ & $9.485867880\cdot10^{-8}$ & $23.329645$ \\
$2.3$ & $35.5865$ & $1.73\cdot10^{-4}$ & $1.124903039\cdot10^{-6}$ & $19.761768$ \\
\bottomrule
\end{tabular}
\end{table}

\begin{proof}
The verifier evaluates weights $1$ through $32$ individually with
\eqref{eq:ec-marker-bound}.  It covers the remaining weights below $k/4$
with contiguous transfer blocks.  From $k/4$ onward, invertibility bounds each
preimage by the output Hamming ball.  The verified sums are the displayed
exact-transfer values.  Theorem~\ref{thm:ec-exact-first-moment} gives the
probability statement.
\end{proof}

The improvement factors are at least $1032$, $181$, and $153$.  The prior
finite calculation assumes an extension of its comparison lemma.  The exact
transfer calculation does not use that assumption.
